\documentclass[12pt]{article}
\usepackage[T2A]{fontenc}
\usepackage[utf8]{inputenc}
\usepackage[russian]{babel}
\usepackage{amsmath, amssymb, amsfonts}
\usepackage{geometry}
\usepackage{fancyhdr}
\usepackage{microtype}
\usepackage{titlesec} % Пакет для управления стилем заголовков

% Настройка Section: по центру, жирный, крупный
\pagestyle{fancy}
% Настройка полей
\geometry{a5paper, margin=1cm, top=2cm, headheight=15pt}

% --- ЖЕСТКАЯ КОРРЕКЦИЯ ЦЕНТРИРОВАНИЯ ---
\pagestyle{fancy}
\fancyhf{} % Полная очистка всех полей (убирает невидимые отступы)
\fancyhead[C]{\thepage} % Номер строго по центру
\renewcommand{\headrulewidth}{0pt} % Убираем линию
\fancyfoot{} % Низ пустой

% Убираем внутренние отступы колонтитулов, которые могут давать перекос
\setlength{\headwidth}{\textwidth} 

% Переопределяем системный стиль plain, чтобы он не сбивал настройки
\makeatletter
\let\ps@plain\ps@fancy
\makeatother

\parindent=0pt
\parskip=0.6em
\titleformat{\section}
  {\normalfont\bfseries\centering}{}{0pt}{}

% Настройка Subsection: отступ 1см, жирный, обычный размер
% {команда}[формат]{стиль}{метка}{отступ_от_метки}{перед_заголовком}[после_заголовка]
\titleformat{\subsection}
  {\normalfont\bfseries}{\hspace{1cm}}{1pt}{}
\titlespacing*{\subsection}{1cm}{2ex plus 1ex minus .2ex}{0.5ex plus .2ex}

% Убираем автоматическую нумерацию, чтобы заголовки были как в оригинале
\setcounter{secnumdepth}{0}


\begin{document}

% --- СТРАНИЦА 5 ---
\setcounter{page}{5}
\section{1.1 Определители}
1. Вычислить $\det \begin{pmatrix} a & b & c \\ b & c & a \\ c & a & b \end{pmatrix}$.

2. Вычислить \\
а) $\det \begin{pmatrix} 0 & -2 & 1 \\ 2 & 0 & 3 \\ -1 & -3 & 0 \end{pmatrix}$; 
б) $\det \begin{pmatrix} -2 & -1 & 1 & 2 \\ -1 & 1 & 2 & -1 \\ 1 & 2 & -2 & -1 \\ 2 & -2 & -1 & 1 \end{pmatrix}$.

3. Проверить, что 
\[ \left( \det \begin{pmatrix} x^1 \\ x^2 \end{pmatrix} \right)^2 = \det \begin{pmatrix} (x^1, x^1) & (x^1, x^2) \\ (x^2, x^1) & (x^2, x^2) \end{pmatrix}, \]
где $x^1 = (x^1_1, x^1_2)$, $x^2 = (x^2_1, x^2_2)$.

4. Вычислить 
\[ \det A_4 = \det \begin{pmatrix} 2 & -1 & 0 & 0 \\ -1 & 2 & -1 & 0 \\ 0 & -1 & 2 & -1 \\ 0 & 0 & -1 & 2 \end{pmatrix}. \]
Чему равен $\det A_n$?

5. Вычислить 
\[ \det V_4 = \det \begin{pmatrix} 1 & 1 & 1 & 1 \\ x_1 & x_2 & x_3 & x_4 \\ x_1^2 & x_2^2 & x_3^2 & x_4^2 \\ x_1^3 & x_2^3 & x_3^3 & x_4^3 \end{pmatrix}. \]
Чему равен $\det V_n$ (определитель Вандермонда)?

6. Доказать, что определитель транспонированной матрицы равен определителю исходной матрицы.

7. Доказать, что определитель кососимметрической матрицы нечетного порядка равен нулю.

8. Доказать, что
\[ \det \begin{pmatrix} 
0 & 1 & 1 & \dots & 1 \\ 1 & 0 & 1 & \dots & 1 \\ 1 & 1 & 0 & \dots & 1 \\ \dots & \dots & \dots & \ddots & \dots \\ 1 & 1 & 1 & \dots & 0 
\end{pmatrix} = (-1)^{n-1}(n-1), \]
где $n$ — порядок определителя.

\newpage

% --- СТРАНИЦА 7 ---
\section{1.2 Прямые и плоскости в $\mathbb{R}^2$ и $\mathbb{R}^3$}

1. Найти уравнение прямой, проходящей через точки $(\xi_1, \xi_2)$;$ (\eta_1, \eta_2)$.

2. Найти уравнение прямой, проходящей через точку $(\xi_1, \xi_2)$ и а) параллельной; б) перпендикулярной прямой $ax_1 + bx_2 + c = 0$.

3. Найти расстояние от точки $(\xi_1, \xi_2)$ до прямой $ax_1 + bx_2 + c = 0$.

4. Найти угол между прямыми $a_1x_1 + b_1x_2 + c_1 = 0$ и $a_2x_1 + b_2x_2 + c_2 = 0$.

5. Найти уравнение срединного перпендикуляра к отрезку, соединяющему точки $(\xi_1, \xi_2)$ и $(\eta_1, \eta_2)$.

6. Найти координаты точки, симметричной точке $(\xi_1, \xi_2)$ относительно прямой $ax_1 + bx_2 + c = 0$.

7. Точки $(\xi_1, \xi_2)$ и $(\eta_1, \eta_2)$ лежат по одну сторону от прямой $ax_1 + bx_2 + c = 0$. Найдите длину кратчайшего пути из одной точки в другую с заходом на данную прямую.

8. Вершины треугольника имеют рациональные координаты. Докажите, что центр его описанной окружности имеет тоже рациональные координаты.

9. Вершины треугольника имеют координаты $(0,0), (a,0)$ и\\ $(b \cos \varphi, b \sin \varphi)$. Найдите координаты а) центра тяжести; б) центра описанной окружности; в) ортоцентра.

10. Докажите, что точки $(x^i_1, x^i_2)$, $i=1,2,3$ лежат на одной прямой тогда и только тогда, когда $\begin{vmatrix} x^1_1 & x^1_2 & 1 \\ x^2_1 & x^2_2 & 1 \\ x^3_1 & x^3_2 & 1 \end{vmatrix} = 0$.

11. Вершины несамопересекающегося многоугольника имеют координаты $(v^i_1, v^i_2)$, $i=1, \dots, n$. Найдите его площадь.

12. Докажите, что квадрат площади параллелограмма, натянутого на векторы $x^1 = (x^1_1, x^1_2)$ и $x^2 = (x^2_1, x^2_2)$, равен $\det \begin{pmatrix} (x^1, x^1) & (x^1, x^2) \\ (x^2, x^1) & (x^2, x^2) \end{pmatrix}$.

13. Пусть $S$ — площадь треугольника со сторонами $a, b, c$. Докажите, что
\[ S^2 = \frac{1}{4} \begin{vmatrix} a^2 & \frac{a^2+b^2-c^2}{2} \\ \frac{a^2+b^2-c^2}{2} & b^2 \end{vmatrix} = -\frac{1}{16} \begin{vmatrix} 0 & a^2 & b^2 & 1 \\ a^2 & 0 & c^2 & 1 \\ b^2 & c^2 & 0 & 1 \\ 1 & 1 & 1 & 0 \end{vmatrix}. \]

14. Найти множество точек плоскости, сумма квадратов расстояний от которых до данных точек $(x^i_1, x^i_2)$, $i=1,\dots,n$ постоянна.

15. Найти уравнение плоскости, проходящей через точки $(0,0,0)$, $(a_1, a_2, a_3)$ и $(b_1, b_2, b_3)$.

16. Найти координаты вектора, параллельного линии пересечения плоскостей $a_1x_1 + a_2x_2 + a_3x_3 + a_4 = 0$ и $b_1x_1 + b_2x_2 + b_3x_3 + b_4 = 0$.

17. Укажите две плоскости, линия пересечения которых параллельна вектору $(a,b,c)$.

18. Найти угол между плоскостями $a_1x_1 + a_2x_2 + a_3x_3 + a_4 = 0$ и $b_1x_1 + b_2x_2 + b_3x_3 + b_4 = 0$.

19. Найти проекцию точки $(u_1, u_2, u_3)$ на плоскость $ax_1 + bx_2 + cx_3 + d = 0$.

20. Найти расстояние от точки $(u_1, u_2, u_3)$ до плоскости $ax_1 + bx_2 + cx_3 + d = 0$.

21. Найти расстояние от точки $(x_1, x_2, x_3)$ до прямой, проходящей через точки $(y_1, y_2, y_3)$ и $(z_1, z_2, z_3)$.

\newpage

% --- СТРАНИЦА 9 ---
22. Докажите, что $\begin{vmatrix} x_1 & x_2 & x_3 \\ y_1 & y_2 & y_3 \\ z_1 & z_2 & z_3 \end{vmatrix}^2 \le (\sum x_i^2)(\sum y_i^2)(\sum z_i^2)$.

23. Пусть $\varphi_{ij}$ — двугранный угол при ребре между $i$-ой и $j$-ой гранями тетраэдра. Докажите, что:
\[ \begin{vmatrix} -1 & \cos \varphi_{12} & \cos \varphi_{13} & \cos \varphi_{14} \\ \cos \varphi_{21} & -1 & \cos \varphi_{23} & \cos \varphi_{24} \\ \cos \varphi_{31} & \cos \varphi_{32} & -1 & \cos \varphi_{34} \\ \cos \varphi_{41} & \cos \varphi_{42} & \cos \varphi_{43} & -1 \end{vmatrix} = 0. \]

24. Пусть $V$ — объем параллелепипеда, натянутого на векторы $v_1$, $v_2$, $v_3$. Докажите, что $V^2 = \det \|(v_i, v_j)\|$.

25. Пусть $V$ — объем тетраэдра $ABCD$; $a=DA, b=DB, c=DC, a_1=BC, b_1=AC, c_1=AB$. Докажите, что
\[ \text{а) } V^2 = \frac{1}{36} \begin{vmatrix} a^2 & \frac{1}{2}(a^2+b^2-c_1^2) & \frac{1}{2}(a^2+c^2-b_1^2) \\ \frac{1}{2}(a^2+b^2-c_1^2) & b^2 & \frac{1}{2}(b^2+c^2-a_1^2) \\ \frac{1}{2}(a^2+c^2-b_1^2) & \frac{1}{2}(b^2+c^2-a_1^2) & c^2 \end{vmatrix}; \]
\[ \text{б) } V^2 = \frac{1}{288} \begin{vmatrix} 0 & c_1^2 & b_1^2 & a^2 & 1 \\ c_1^2 & 0 & a_1^2 & b^2 & 1 \\ b_1^2 & a_1^2 & 0 & c^2 & 1 \\ a^2 & b^2 & c^2 & 0 & 1 \\ 1 & 1 & 1 & 1 & 0 \end{vmatrix}. \]

26. Пусть $R$ — радиус описанной окружности треугольника $ABC$. Докажите, что
\[ 2R^2 \begin{vmatrix} 0 & c^2 & b^2 & 1 \\ c^2 & 0 & a^2 & 1 \\ b^2 & a^2 & 0 & 1 \\ 1 & 1 & 1 & 0 \end{vmatrix} + \begin{vmatrix} 0 & c^2 & b^2 \\ c^2 & 0 & a^2 \\ b^2 & a^2 & 0 \end{vmatrix} = 0. \]

\newpage

% --- СТРАНИЦА 10 ---
\section{1.3 Подпространства и аффинные многообразия в $\mathbb{R}^n$}

1. Найти основание $\eta$ перпендикуляра, опущенного из точки $\xi$ на гиперплоскость $(x, a) = 0$, где $|a|=1$.

2. Найти уравнение гиперплоскости, проходящей через точку $\xi$ параллельно гиперплоскости $(x, a) = b$.

3. Найти расстояние от точки $\xi$ до гиперплоскости $(x, a) = 0$, где $|a|=1$.

4. Найти уравнение гиперплоскости, проходящей через точки\\ $(0,\dots,0)$, $(x^k_1, \dots, x^k_n)$, где $k=1,\dots,n-1$.

5. Описать систему уравнений, задающих подпространство, по которому пересекаются гиперплоскости $\sum a^j_i x_i = 0, j=1,\dots,k$.

6. Найти угол $\varphi$ между $(n-1)$-мерными гранями правильного $n$-мерного симплекса.

7. Представьте гиперплоскость $x_1 + x_2 + \dots + x_n = 0$ как линейный образ $\mathbb{R}^{n-1}$.

8. Представьте прямую $x_i = t, 1 \le i \le n$, как пересечение гиперплоскостей.

9. Докажите, что
\[ \left| \det \begin{pmatrix} x^1 \\ \vdots \\ x^n \end{pmatrix} \right| \le \prod_{i=1}^n |x^i|, \quad x^i = (x^i_1, \dots, x^i_n), \ 1 \le i \le n \]
(неравенство Адамара).

\newpage

\section{1.4 Теория линейных уравнений и геометрия.}

В этом параграфе $\mathbb{R}^n$ — $n$-мерное пространство векторов 
\[ \{x=(x_1, \dots, x_n) \mid x_i \in \mathbb{R}\} \]
(без евклидовой структуры).

Подпространство $\mathbb{R}^n$ - это множество в $\mathbb{R}^n$ такое, что сумма векторов и произведение вектора на число не выводят из данного множества, аффинное множество - сдвиг подпространства.

1. Доказать, что любое подпространство в $\mathbb{R}^n$ замкнуто.

2. Доказать, что для того, чтобы множество в $\mathbb{R}^n$ было подпространством, необходимо и достаточно, чтобы оно было решением конечной системы однородных равенств.

3. Доказать, что любое подпространство допускает двойное описание: оно есть образ некоторого $\mathbb{R}^k$ при линейном отображении с одной стороны, а с другой — оно есть решение некоторой системы однородных равенств.

4. Доказать, что если любые $k$ гиперплоскостей семейства пересекаются, то пересечение всех гиперплоскостей есть аффинное множество размерности $\ge n-k$.

5. Доказать, что аффинное многообразие размерности $k$ в $\mathbb{R}^n$ допускает двойное описание: оно есть образ $\mathbb{R}^k$ при аффинном отображении (т.е. при отображении $x \rightarrow Ax + b$, где $A:\mathbb{R}^k \to \mathbb{R}^n$ - линейное отображение, а $b \in \mathbb{R}^n$), а с другой стороны оно есть решение некоторой системы равенств.

6. Доказать, что размерность подпространства (аффинного множества) равна наименьшему из тех $k$, при которых оно есть линейный (аффинный) образ $\mathbb{R}^k$.

7. Доказать эквивалентность двух определений $k$-мерного подпространства в $\mathbb{R}^n$:
\begin{itemize}
    \item[а)] оно есть образ $\mathbb{R}^k$ при линейном отображении и не является таковым для $\mathbb{R}^l$, $l < k$;
    \item[б)] оно есть пересечение $n-k$ гиперплоскостей, а любое пересечение меньшего числа гиперплоскостей содержит его собственное подмножество.
\end{itemize}

8. Пусть $A = (a_{ij})$ — матрица размеров $n \times m$. Доказать, что следующие определения приводят к одному и тому же числу (называемому рангом матрицы $A$):
\begin{itemize}
    \item[а)] ранг матрицы $A$ есть размерность подпространства $A\mathbb{R}^n$ в $\mathbb{R}^m$;
    \item[б)] ранг матрицы $A$ есть максимальный порядок отличного от нуля минора;
    \item[в)] ранг матрицы $A$ есть максимальное число линейно независимых столбцов (строк) матрицы $A$.
\end{itemize}

9. Для того, чтобы система $Ax = b$ была разрешима, необходимо и достаточно, чтобы ранг матрицы $A = (a_{ij})$ совпадал с рангом расширенной матрицы 
\[
\overline{A} = \begin{pmatrix}
a_{11} & \cdots & a_{1n} & b_1 \\
\vdots & \ddots & \vdots & \vdots \\
a_{m1} & \cdots & a_{mn} & b_m
\end{pmatrix}.
\]

\section{1.5. Ответы, указания, решения}

\subsection{1.2. Прямые и плоскости в $\mathbb{R}^2$ и $\mathbb{R}^3$}

2. \textit{Ответ:} а) $ax_1 + bx_2 = a\xi_1 + b\xi_2$; б) $-bx_1 + ax_2 = -b\xi_1 + a\xi_2$.

3. \textit{Ответ:} $\frac{|a\xi_1 + b\xi_2 + c|}{\sqrt{a^2 + b^2}}$.

4. \textit{Ответ:} $\cos \varphi = \frac{a_1a_2 + b_1b_2}{\sqrt{a_1^2 + b_1^2} \sqrt{a_2^2 + b_2^2}}$.

5. \textit{Ответ:} $(x_1 - \xi_1)^2 + (x_2 - \xi_2)^2 = (x_1 - \eta_1)^2 + (x_2 - \eta_2)^2$, т.е.

$2x_1(\xi_1 - \eta_1) + 2x_2(\xi_2 - \eta_2) = \xi_1^2 - \eta_1^2 + \xi_2^2 - \eta_2^2$.

6. \textit{Ответ:} $\left( \xi_1 - \frac{2a(a\xi_1 + b\xi_2 + c)}{a^2 + b^2}, \xi_2 - \frac{2b(a\xi_1 + b\xi_2 + c)}{a^2 + b^2} \right)$.

7. \textit{Ответ:} в случае прямой $x_1 = x_2$ длина равна $\sqrt{(\xi_1 - \eta_2)^2 + (\xi_2 - \eta_1)^2}$.

8. \textit{Указание:} воспользуйтесь задачей 5.

9. \textit{Ответ:} а) $\left( \frac{b \cos \varphi + a}{3}, \frac{b \sin \varphi}{3} \right)$; б) $\left( \frac{a}{2}, \frac{b - a \cos \varphi}{2 \sin \varphi} \right)$; в) $(b \cos \varphi, (a - b \cos \varphi) \text{ctg} \varphi)$.

11. \textit{Ответ:} $\frac{1}{2} \sum\limits_{k=1}^n \begin{vmatrix} v^k_1 & v^{k+1}_1 \\ v^k_2 & v^{k+1}_2 \end{vmatrix}$.

\textit{Указание:} воспользуйтесь формулой для площади треугольника, одна вершина которого лежит в начале координат.

12. \textit{Указание:} рассмотрите произведение матриц $\begin{pmatrix} x^1_1 & x^1_2 \\ x^2_1 & x^2_2 \end{pmatrix} \begin{pmatrix} x^1_1 & x^2_1 \\ x^1_2 & x^2_2 \end{pmatrix}$.

13. \textit{Указание:} Воспользуйтесь задачей 13. Для доказательства второго равенства вычтите первую строку определителя $4 \times 4$ из всех остальных строк, кроме последней, а затем вычтите первый столбец из всех остальных столбцов, кроме последнего.

14. \textit{Ответ:} окружность $\sum (x_1 - x^i_1)^2 + \sum (x_2 - x^i_2)^2 = d$.

15. \textit{Ответ:} $\begin{vmatrix} a_2 & a_3 \\ b_2 & b_3 \end{vmatrix} x_1 - \begin{vmatrix} a_1 & a_3 \\ b_1 & b_3 \end{vmatrix} x_2 + \begin{vmatrix} a_1 & a_2 \\ b_1 & b_2 \end{vmatrix} x_3 = 0$.

\textit{Указание:} $\begin{vmatrix} x_1 & x_2 & x_3 \\ a_1 & a_2 & a_3 \\ b_1 & b_2 & b_3 \end{vmatrix} = 0$ при $x = a$ и $x = b$.

16. \textit{Ответ:} $\left( \begin{vmatrix} a_2 & a_3 \\ b_2 & b_3 \end{vmatrix}, -\begin{vmatrix} a_1 & a_3 \\ b_1 & b_3 \end{vmatrix}, -\begin{vmatrix} a_1 & a_2 \\ b_1 & b_2 \end{vmatrix} \right)$.

17. \textit{Указание:} при $c \ne 0$ плоскость $\alpha x_1 + \beta x_2 + \gamma x_3 = 0$, где $\gamma = -(\alpha a + \beta b)/c$, параллельна вектору $(a, b, c)$.

\newpage
18. \textit{Указание:} вектор $(a, b, c)$ перпендикулярен плоскости 

$ax_1 + bx_2 + cx_3 + d = 0$.

19. \textit{Ответ:} $(u_1 + \lambda a, u_2 + \lambda b, u_3 + \lambda c)$, где $\lambda = -\frac{au_1 + bu_2 + cu_3 + d}{a^2 + b^2 + c^2}$.

20. \textit{Ответ:} $\frac{|au_1 + bu_2 + cu_3 + d|}{\sqrt{a^2 + b^2 + c^2}}$.

21. \textit{Ответ:} $d^2 = \frac{|u|^2|v|^2 - \langle u, v \rangle^2}{|u|^2 + |v|^2 - 2\langle u, v \rangle}$, где $u$ и $v$ — векторы, идущие из точки $(x_1, x_2, x_3)$ в точки $(y_1, y_2, y_3)$ и $(z_1, z_2, z_3)$.

22. \textit{Указание:} объем параллелепипеда не превосходит произведения длин его ребер.

23. \textit{Первое решение.} Система уравнений 
\[ \begin{cases} -x_1 + x_2 \cos \varphi_{12} + x_3 \cos \varphi_{13} + x_4 \cos \varphi_{14} = 0 \\ x_1 \cos \varphi_{21} - x_2 + x_3 \cos \varphi_{23} + x_4 \cos \varphi_{24} = 0 \\ x_1 \cos \varphi_{31} + x_2 \cos \varphi_{32} - x_3 + x_4 \cos \varphi_{34} = 0 \\ x_1 \cos \varphi_{41} + x_2 \cos \varphi_{42} + x_3 \cos \varphi_{43} - x_4 = 0 \end{cases} \]
имеет ненулевое решение $x_i = S_i$, где $S_i$ — площадь $i$-ой грани.

\textit{Второе решение.} Пусть $n^i$ — вектор внешней нормали к $i$-ой грани. Тогда $\langle n^i, n^j \rangle = -\cos \varphi_{ij}$ и $\det \| \langle n^i, n^j \rangle \| = 0$, так как это есть определитель Грама порядка 4 в 3-мерном пространстве.

26. \textit{Указание:} Объем тетраэдра $ABCO$, где $O$ — центр описанной окружности, равен 0. В определителе $\begin{vmatrix} 0 & c^2 & b^2 & R^2 & 1 \\ c^2 & 0 & a^2 & R^2 & 1 \\ b^2 & a^2 & 0 & R^2 & 1 \\ R^2 & R^2 & R^2 & 0 & 1 \\ 1 & 1 & 1 & 1 & 0 \end{vmatrix}$ прибавьте последние строку и столбец, умноженные на $-R^2$, к предпоследним строке и предпоследнему столбцу. Полученный определитель разложите по предпоследней строке.

\subsection{1.3. Подпространства и аффинные многообразия в $\mathbb{R}^n$}

1. \textit{Ответ:} $\eta = \xi - \langle \xi, a \rangle a$.

2. \textit{Ответ:} $\langle x, a \rangle = \langle \xi, a \rangle$.

3. \textit{Ответ:} $|\langle \xi, a \rangle|$.

4. \textit{Ответ:} $\sum a_i x_i = 0$, где $a_i = (-1)^i \det \begin{pmatrix} x^1_1 & \dots & x^1_i & \dots & x^1_n \\ \dots & \dots & \dots & \dots & \dots \\ x^{n-1}_1 & \dots & x^{n-1}_i & \dots & x^{n-1}_n \end{pmatrix}$.

5. \textit{Ответ:} $\sum b_i x_i = 0$, где $b_i = (-1)^i \det \begin{pmatrix} a^1_1 & \dots & a^1_i & \dots & a^1_n \\ \dots & \dots & \dots & \dots & \dots \\ a^k_1 & \dots & a^k_i & \dots & a^k_n \\ t^{k+1}_1 & \dots & t^{k+1}_i & \dots & t^{k+1}_n \\ \dots & \dots & \dots & \dots & \dots \\ t^{n-1}_1 & \dots & t^{n-1}_i & \dots & t^{n-1}_n \end{pmatrix}$ \\
($t^j_i$ — произвольные параметры).

6. \textit{Ответ:} $\cos \varphi = \frac{1}{n}$.

\textit{Указание:} Пусть $S$ — $(n-1)$-мерный объем грани симплекса. Тогда $(n-1)$-мерный объем проекции каждой из $n$ боковых граней на основание равен $S \cos \varphi$, а их сумма равна $S$.

\end{document}
